\documentclass[11pt, a4paper]{common/document_template}
\usepackage{common}

\begin{document}

\begin{titlepage}
  \begin{center}
	\Huge{Задачи по комбинаторике}
  \end{center}
\end{titlepage}

\begin{exercise}
Вокруг стола расселись рыцари и лжецы. Каждый из них сказал о своем соседе справа, правдив тот или лжив. Известно, что на основании этих заявлений можно однозначно определить, какую долю от присутствующих составляют рыцари. Чему она равна?
\end{exercise}

\end{document} 
